% WORKING DRAFT: PRIMARY RESULTS REMAIN V2 UNTIL THE V3 COMPLETION GATE.
\section{Evaluation}
\label{sec:evaluation}

We evaluate failure-induced workflow delay, the incremental value of Peer HBM
over Peer DRAM, and the conditions under which workflow Slack changes
protection. Discrete-event simulations measure modeled completion time;
separate native-selector observations test selection without assigning it a
serving-time benefit. Negative and null results constrain the design's claims.

\subsection{Methodology}
\label{sec:eval-method}

\paragraph{Simulator and timing model.}
MoonKVSim models dependency release, queues, asynchronous prediction,
checkpoint publication and recovery on one event clock. Transfers contend for
rank H2D/D2H, host DRAM/NIC and fabric resources; pending copies consume capacity
and every TP shard must commit before use. Worker failures remove a TP group;
host failures also remove DRAM. Audits check token counts, fencing, capacity
and release, not generated-token identities.
Pinned AIConfigurator HYBRID estimates\cite{eval2026aiconfigurator} provide
dense Qwen3-8B TP1 and Qwen3-32B TP4 A100-SXM durations through bounded
interpolation over 196 shapes; decode sums individual steps. Mixed
empirical/silicon provenance is retained, without end-to-end calibration.
Each grid has 49 prefixes (0--98,304, step 2,048), prefill extensions of
1/64/256 tokens, and one-token decode.
FCFS permits one active invocation per worker, excluding continuous batching
and production-throughput conclusions.

\paragraph{Workload and split.}
Constructed, evidence-provided FanOutQA replays\cite{eval2024fanoutqa} retain
explicit dependencies. A frozen task-ID order/context guard selected 160 of
310 development questions; retaining workflows whose replayed request contexts
all fit 40,960 tokens leaves 74. A predeclared workflow-ID hash gives 12 training workflows (95 tasks) and
62 evaluation workflows (503 tasks). These exclusions constrain
representativeness. Each DAG runs once; each task is a distinct,
single-invocation Agent. Pinned tokenizers encode prompts and gold answers;
reservation uses a declared 512-token output limit, not eventual output length
or future calls.
Separate constructed four-Agent fork--joins test branch skew, repeated-Agent
history, long decode with synthetic demand, and shared hosts: eight workflows,
four TP1 workers, seeds 7/19/41 varying token identities/length jitter, and
16,384-token output caps. Shared-host cases colocate two workers per host.
These are interventions, not natural multi-call traces.

\paragraph{Resources and prediction.}
Primary runs use eight independently hosted workers (8/32 GPUs for 8B/32B).
Each GPU has 80\,GiB minus profile allocations and a 4\,GiB runtime allowance;
each host has 256\,GiB Peer DRAM. Effective H2D/D2H, NIC/fabric and DRAM rates
are 32/25/100\,GB/s. Per-rank Peer HBM starts at 128\,MiB, grows by 64\,MiB
after stable windows, and caps at 4\,GiB with 2\,GiB foreground reserve.
Blocks/checkpoints span 64/256 tokens; prefill chunks span 256.
We use $K=16$, $\beta=0.9$, $\epsilon=0.1$, $\lambda=10$\,ms and
16 selector-state and eight counter bytes per block/layer/rank.

A submission-time hashed-text regressor with 40 Context features predicts
same-Agent final completion minus submission, including queueing. Six healthy
training runs (two models, three loads) yield 570 labels from 12 workflows.
Grouped leave-one-workflow-out calibration holds all three loads together and
fits role priors only on other workflows. The guard is the training 90th
percentile of workflow-maximum positive Slack overestimation: 0.910/0.995\,s
for 8B/32B, referenced to healthy realized completion, not queue-free timing.
All training Agents have one invocation; validation therefore does not test
future-call prediction. This empirical guard is not a probabilistic delay bound.
Evaluation labels tune neither predictor nor guard; semantic-encoder accuracy
is unvalidated. Reconstruction estimates share execution's profile, an idealization.
Checkpoint-aligned training-median context (8,704 tokens), dense demand/metadata and uniform
layer compute fix $L=11/1$ offline; constructed controls use declared priors
and $L=4$. Earlier resource/Slack controls retain a service-only proxy using
prompt work and training role-median outputs, excluding queueing.

\paragraph{Policies and faults.}
One runtime implements Restart; All-DRAM (no exposure threshold, same ordering);
DRAM (workflow selection without extra HBM); Static-HBM (128\,MiB initially,
no growth, foreground shrink allowed); and \system. These controlled variants
are not prior-system reproductions; DRAM isolates the HBM increment.
Arrivals use training work, eight workers and nominal
$\rho\in\{0.3,0.65,0.9\}$, not measured utilization.
Each model/load/policy has F0 and five faults (180 runs), targeting equally
spaced source-order workflows, alternating sink/middle Agents at 80\% of
invocation-total logical tokens: four prefill faults and one decode fault.
The 135 constructed runs instead target 80\% of each phase.
Detection takes 1\,ms; failed capacity remains lost. The post-failure work-rate
proxy $8\rho/7$ reaches 1.03. Matched logical identities/positions can differ
in wall time and co-victims; splitting prefill can also change profiled cost.

\paragraph{Endpoints and provenance.}
WCT spans external arrival to last-Agent completion. Co-primary endpoints
are same-workflow $\Delta\mathrm{WCT}=\mathrm{WCT}_{F1}-\mathrm{WCT}_{F0}$
and interrupted-Agent failure-to-first-new-token latency, decomposed into
detection, retry queue and post-admission service. Each F1 subtracts its own
policy's F0 for that workflow. Sink/non-sink identities, all-workflow means,
and adverse outcomes remain separate; cohort outcomes can oppose target gains.
Five fixed faults
are descriptive; workflows share queues, reused F0 is no new sample, and
62-workflow interpolated p95/p99 are not reliable population tails.
Deterministic reruns receive no confidence intervals.
Artifacts seal sources, inputs, predictors/guards and accounting. The original
483 conditions comprise 180 service-proxy primary, 135 structural, 126 sensitivity,
32 capacity/Slack and ten post-hoc cases; new primary runs change only predictor,
priors and guard. DMA/quota extensions add 24/72 conditions, not independent
repetitions; DMA retains the service proxy, quota uses the new predictor.
Artifacts retain absolute F0/F1 WCT, checkpoint lag, transfer bytes,
occupancy/byte-time, active co-victims and queued-history restores separately,
with full transfer intervals for selected runs.
Common runtime checks cover causal KV/local tails, queued-history
restoration, accepted-EOS cancellation and 74 regressions; historical versions
remain separate.

\subsection{Primary Workflow Replay}
\label{sec:eval-primary}
\paragraph{Prediction check.}
The six healthy held-out trajectories give mean absolute errors of
2.027/2.201\,s for 8B/32B versus role-prior baselines of 2.073/2.248\,s.
At $\rho=0.9$, error rises to 3.290/3.384\,s with substantial underprediction.
Each Agent receives identical Contexts and predictions across loads despite
different queueing; the predictor has no load/queue feature. Healthy errors
do not validate failure-conditioned predictions.

% TODO: Replace every following primary number/figure/table and recheck attribution
% against the completed v3 artifacts. Until then this subsection remains v2.
Figure~\ref{fig:eval-primary} retains every sink/non-sink target.
For 8B, DRAM reduces mean target penalty from 0.530 to 0.062\,s at
$\rho=0.3/0.65$, but worsens 1.994 to 2.070\,s at $\rho=0.9$.
Cohort means move oppositely: 0.744 to 0.766\,s at $\rho=0.65$ and
3.956 to 3.587\,s at $\rho=0.9$.
For 32B at $\rho=0.3$, mean target penalty improves 1.427 to 1.124\,s,
yet the worst target worsens 3.876 to 5.299\,s: changing retry workers
adds 1.453\,s queueing while saving 29\,ms service.

Retry queue contributes 62.22--90.63\% of DRAM's mean Agent latency
(Table~\ref{tab:agent-failure-latency}); retries cannot preempt busy Peers.
Three sink targets fail at 257--356 tokens with a 256-token checkpoint,
limiting avoidable reconstruction. Per-identity decompositions remain in the artifact.

The target-workflow HBM increment ranges from $-0.0026$ to 1.442\,ms for 8B
and from 0 to 0.137\,ms for 32B across the 15 model-specific fault conditions.
Static-HBM and \system have identical individual WCT, so these runs do not
demonstrate an adaptation advantage. Healthy mean WCT (6.23--16.65\,s)
is also identical across policies despite nonzero traffic, reflecting the
model's execution and interference assumptions.

\begin{table}[t]
\centering
\setlength{\tabcolsep}{3pt}
\renewcommand{\arraystretch}{1.08}
\caption{Mean failure-to-first-new-token latency over five fixed identities.
HBM gain is DRAM minus AegisKV; positive is better. No confidence intervals.}
\label{tab:agent-failure-latency}
\begin{tabular}{@{}lrrrr@{}}
\hline
Model & $\rho$ & Restart & DRAM & HBM gain \\
 & & (s) & (s) & (ms) \\
\hline
Qwen3-8B & 0.3 & 0.956 & 0.488 & 0.220 \\
Qwen3-8B & 0.65 & 0.956 & 0.488 & 0.219 \\
Qwen3-8B & 0.9 & 2.694 & 1.970 & 0.508 \\
Qwen3-32B & 0.3 & 1.910 & 1.607 & 0.027 \\
Qwen3-32B & 0.65 & 1.144 & 0.551 & 0.027 \\
Qwen3-32B & 0.9 & 2.222 & 1.253 & 0.036 \\
\hline
\end{tabular}
\end{table}

\begin{figure*}[t]
  \centering
  \includegraphics[width=\textwidth]{figures/eval-primary-wct.pdf}
  \caption{Target-workflow penalty for five fixed identities at each model/load.
  S/N denote sink/non-sink Agents. Each marker subtracts own-policy healthy WCT
  for the same workflow; offsets separate coincident policies. Logical targets
  match across policies; negative penalties remain visible. No confidence
  intervals. Nominal load is an arrival-generation parameter.}
  \label{fig:eval-primary}
  \Description{Six panels compare target-workflow paired delay at five fixed sink or non-sink failure identities for two models and three nominal loads. Five policies use offset markers; negative penalties are retained.}
\end{figure*}

Against All-DRAM, the exposure threshold saves 10.21/5.70/$-0.095$\% of
healthy backup traffic at increasing 8B loads and 10.81/7.98/1.16\% for 32B,
including cancellation bytes. Cohort means match despite a $\pm4.113$\,s
workflow swap in one 32B high-load fault. In the two logged \system F1
runs at $\rho=0.65$, fallback covers 3,122/5,586 and 3,605/6,518 unfinished-WF
control-window observations (8B/32B), all from nonpositive prediction minus
submission age. These are event counts, not time fractions or healthy-run rates.

In the second 8B high-load fault, HBM advances the target by 1.10\,ms,
moving a completion across a fork submission and changing routing counts.
Another workflow finishes 3.41\,s later; eight gain 1.10\,ms and 53 are
unchanged, increasing cohort WCT sum by 3.40\,s. Fault time, worker and
co-victims match. Recorded submissions/completions and worker choices support
this reversal; local $G_a$ does not guarantee global WCT improvement.

\subsection{Workflow Structure and Recovery}
\label{sec:eval-structure}
Across the 135 constructed runs, DRAM worsens six of ten scenario/phase
mean target penalties and two cohort means (Figure~\ref{fig:eval-structure}).
Both worsen for skewed prefill (14.43/2.89\,s, target/cohort) and long-decode
decode (21.24/6.63\,s); the latter's seed 19 worsens by 49.70/15.34\,s.
Long-decode prefill also worsens the target by 30.62\,s. Extra HBM changes
the ten mean target penalties by less than 0.875\,ms, and by below one
microsecond in long-decode cases. There is no uniform checkpoint improvement
or substantial incremental HBM benefit.

\begin{figure*}[t]
  \centering
  \includegraphics[width=\textwidth]{figures/eval-structural-wct.pdf}
  \caption{Target-workflow penalty in five constructed four-Agent workloads,
  eight workflows per run. Gray segments connect matched seeds; colored ranges
  show min/max and horizontal marks means. Each F1 subtracts own-policy F0.
  Agent, second invocation and logical position match across policies;
  prefill/decode positions differ. Negative penalties and overlapping
  DRAM/AegisKV points remain visible.}
  \label{fig:eval-structure}
  \Description{Two panels show target-workflow delay for prefill and decode faults across five constructed scenarios, with three matched seed points per policy. Negative penalties remain visible. Shared-host failures have large variation; checkpointing can worsen the target workflow.}
\end{figure*}

Endpoint choice can reverse the conclusion. In shared-host prefill/decode
cases, DRAM worsens the mean target penalty by 30.00/41.01\,s while reducing
the all-workflow mean penalty by 16.11/19.01\,s. History-reuse prefill likewise
worsens the target by 12.05\,s but improves the cohort mean by 3.33\,s.
Permanent capacity loss, queue placement, and completion-dependent routing
redistribute delay across workflows; a local gain $G_a$ is not a global WCT
guarantee. The audit separately records 108 active interrupted recoveries and
279 queued-history restorations, without counting the latter as active
co-victims. All per-seed target/cohort endpoints and recovery decompositions
remain in the artifact.

The long-decode control's two synthetic four-block banks both fit $K=16$;
it cannot establish native hotspot robustness under a binding budget.
The 32/128-workflow CPU/RSS controls use a historical runtime and establish
neither current recovery results nor physical controller overhead.

\subsection{Operating Range and Controls}
\label{sec:eval-sensitivity}
The 126 one-factor controls use four complete held-out workflows (52 Agents),
fixed arrivals/priors, and the longest-prompt Agent in the third workflow:
20,283 prompt tokens, interrupted at token 16,228. Input length selects this
target before outcomes. Each condition subtracts its own F0; shared center
runs are not additional samples. This is one fixed fault, not a population.

In Figure~\ref{fig:eval-resources}, reducing H2D from 64 to 8\,GB/s raises the
DRAM target penalty from 24.02 to 283.95\,ms; HBM saves another 0.72/5.78\,ms.
At 32\,GB/s, restart/DRAM/\system penalties are 1.8023/0.0610/0.0595\,s,
with a 1.44\,ms HBM increment. At 0.25/1\,GiB Peer DRAM, the target loses its
checkpoint and both backup policies revert to restart; peak host occupancy
reaches 98.45\%. This is a capacity boundary, not unconstrained DRAM operation.
A zero/32\,MiB HBM cap gives the same target penalty as DRAM.

A 32-run control crosses 2/4/8/256\,GiB DRAM with DRAM, All-DRAM,
\system and zero DAG Slack. Current and zero Slack have identical individual
WCT; Current saves 4.58--4.66\% healthy backup bytes, including cancellation.
The target checkpoint needs 2.215\,GiB, cannot fit at 2\,GiB, and stays at
16,128 tokens for every variant at larger budgets: traffic savings without
competition-induced WCT gains.

A post-hoc ten-run fork--join diagnostic declares 0/80\,s branch Slack
without Peer HBM. A 1.25\,GiB peer budget fits either checkpoint but not both.
DAG Slack retains the critical branch's 1,280-token checkpoint; zero Slack
keeps the noncritical branch's 8,448 tokens. At the same fault with no retry
queue, penalties are 23.97/134.79\,ms. At 4\,GiB both yield 23.97\,ms.
This is a conditional selection witness, not natural prediction evidence.

\begin{figure*}[t]
  \centering
  \includegraphics[width=\textwidth]{figures/eval-resource-sensitivity.pdf}
  \caption{One fixed prefill fault on four complete FanOutQA workflows,
  Qwen3-8B TP1, eight workers. Points are target-workflow F1 minus own-condition
  F0 WCT; no confidence interval is implied. $L=11$, $K=16$, and initial HBM
  quota $\min(128\,\mathrm{MiB},\mathrm{cap})$. Delay, bandwidth and DRAM axes
  are logarithmic; the HBM axis is linear from 0 to 32\,MiB and logarithmic
  thereafter. DRAM-capacity markers show only tested settings; no capacity
  threshold is interpolated.}
  \label{fig:eval-resources}
  \Description{Three panels vary H2D bandwidth, host DRAM and per-rank HBM. Faster H2D lowers recovery delay. Small DRAM budgets lose the target checkpoint; low HBM quotas match DRAM-only.}
\end{figure*}

Prediction multipliers 0.5/1/2, $\lambda=0/0.01/1$\,s and predictor latency
0.1--200\,ms leave target penalty unchanged; at 200\,ms all 52 replies
exceed the 100\,ms TTL. These check non-crossing code paths, not accuracy or
robustness. Detection controls use 1/10/100\,ms.

An additional 24-run extension tests lower DMA slowdown coefficients
0.002/0.005/0.01/0.02 with the same service predictor and fixed fault.
\system's healthy four-workflow WCT sum increases by 0.045--0.447\,s, while
its absolute F1 sum improves by 0.365--0.816\,s against restart.
At coefficients 0.05/0.1, healthy mean WCT
increases by 3.33/6.66\%, from 8.398\,s to 8.678/8.957\,s;
even F1 cohort WCT sums worsen by 0.39/0.74\,s.
For a hypothetical mixture of F0 and this exact
F1 with probability $p$, advantage is $pD-H$: healthy cost $H$ and reduction
in paired fault penalty $D$ must use the same objective. The four-workflow
break-even thresholds at the four lower coefficients are
0.052/0.131/0.266/0.551. At 0.1 they are 0.141 for the target and 1.494
for the four-workflow sum, so no feasible $p\in[0,1]$ offsets cohort cost.
Neither these hypothetical probabilities nor the uncalibrated coefficients
estimate production behavior; no continuous crossover is inferred.
Finally, $Z=64$ also limits forwards to 64 tokens, whereas $Z=1024$ retains
256-token chunks. The former raises even restart's penalty to 17.71\,s:
a coupled compute/checkpoint intervention. Full controls and adverse outcomes
remain in the artifact.

\begin{figure*}[t]
  \centering
  \includegraphics[width=\textwidth]{figures/eval-native-component.pdf}
  \caption{Raw-KV component differences at 32\,GB/s, $K=16,L=4$.
  Positive means Current is faster. Each document first averages checkpoints
  32,000/32,256; markers then show the three-document mean and ranges their
  min/max. Compute settings are categorical assumptions, not measurements.
  The frequency oracle is not latency-optimal. Prepared residency and auxiliary
  state are assumed; these component times are not native-kernel or WCT results.}
  \label{fig:eval-native-component}
  \Description{Three model panels show signed raw-KV component time differences against Current for DRAM, Recency and a future-frequency reference. Current is faster than DRAM but can be slower than Recency and the reference.}
\end{figure*}

\paragraph{Quota pressure.}
A separate 72-run diagnostic uses the same four workflows with the trained
wall-time predictor, priors and guard. The allocatable pool equals the maximum
declared foreground allocation plus a 2\,GiB reserve and 0/32/128\,MiB margin;
every invocation fits alone. DMA factors 0/0.005 and F0/F1 cross restart,
DRAM, fixed 32/128/4096\,MiB targets and adaptive quotas. Fixed targets recover
each window within the live safety ceiling, unlike the earlier permanently
clamped Static-HBM. Pressure occurs in 12/36 F0 runs and no F1 runs, with no
capacity-denied admissions. Both 32-to-24\,MiB adaptive decreases occur after
the target workflow finishes. Adaptive and fixed 128/4096\,MiB have identical
individual workflow endpoints in all 24 matched comparisons. Their 1.444\,ms
target advantage over DRAM cannot be attributed to AIMD. This constructed
small-pool boundary does not demonstrate recovery under pressure or production
HBM utilization.

\subsection{Native Selection and its Limits}
\label{sec:eval-native}
We separately replay sealed NOSA-8B, NSA-8B, and MiniCPM4.1-8B selectors on the
same three PG19 documents. We union blocks across KV groups per query, start
counters at query zero, and apply decay, hysteresis, and the earlier-block tie
rule every 256 tokens. All 13,824 checked sets match the deployed ranking
function. This assumes continuous protection and immediate checkpoint visibility;
it does not reproduce publication gaps or native chunked-prefill execution.



For NOSA/NSA/MiniCPM respectively, $K=16,L=4$ gives Current distinct coverage
of 8.31\%, 4.20\%, and 4.06\%. Repeated query--block event coverage is
15.70\%, 11.46\%, and 5.37\%, versus Recency's 15.18\%, 15.44\%, and 11.45\%
and the oracle's 15.83\%, 16.15\%, and 11.45\%. These observations do not
support ``few misses'' at this budget. They also differ from the dense WCT
models and $L=11/1$ configurations.

At $K=64$, Current's distinct-coverage advantage over Recency is
2.36, 0.97, and $-1.07$ percentage points (NOSA/NSA/MiniCPM);
MiniCPM also loses on event coverage. Full budget curves remain in the artifact.
These denominators are not runtime cache-hit rates or WCT benefits.

An 864-condition raw-KV component (Figure~\ref{fig:eval-native-component})
uses the same three books per model,
two adjacent checkpoints, four residency choices, three H2D rates and four
declared compute times. Only prepared $K=16,L=4$ residency differs;
background order is shared. Each book first averages its two checkpoints.
At 32\,GB/s and zero compute, Current saves 131.072/131.072/128.683\,$\mu$s
over DRAM (NOSA/NSA/MiniCPM), but matches Recency on the first two and
trails it by 2.389\,$\mu$s on MiniCPM. At 256\,$\mu$s assumed layer compute,
it trails Recency by 209.436/218.124\,$\mu$s on NOSA/NSA.
Prepared residency helps without establishing Current's superiority.
This component assumes layer-entry demand release, local tails and prepared
auxiliary state, excluding preparation cost. Per-object readiness differs
from the core's atomic LoadBatch; the future-frequency reference is not a
latency bound. Uncalibrated compute and native full-prefill traces establish
neither native-kernel recovery timing, chunked-resume equivalence nor WCT.

\subsection{Scope of the Evidence}
\label{sec:eval-limits}
Unmodeled CPU/kernel costs and batching, and uncalibrated DMA interference,
preclude physical-overhead claims. Immediate continuations exercise queued
history recovery without tool/think idle gaps.

For one Agent's dense checkpoint, HBM holds at most $LK$ historical blocks per
rank; other demanded blocks need H2D. This partial residency does not exercise
Equation~\ref{eq:depth}'s no-wait first-$L$-layers premise. Bytes/bandwidth
explains local service savings, not global WCT under rerouting. FCFS also
does not establish quota adaptation's benefit under changing batch memory pressure.

Asynchronous publication can leave more than $Z$ uncommitted tokens.
A historical, pre-freeze linear-cost control accepts 6,247 tokens with
committed frontier 1,024: a 5,223-token tail despite $Z=256$.
Likewise, protection tolerance and modeled depth are parameters, not
per-workflow delay or transfer-hiding guarantees.
